\section{Introduction}
\label{sec:intro}

An LLM agent deployed over days or weeks accumulates an interaction
history far larger than any context window. It cannot attend to
everything, so on every consolidation step it makes a triage decision:
which observations to \emph{encode} (and how deeply), which to
\emph{forget}, and which to \emph{retrieve} when a new task arrives.
This decision is not cosmetic --- it determines whether the agent
remembers a user's stated preference three weeks later, whether it
recalls the one reliable fact among a hundred distractors, and whether
its working memory is dominated by recent chatter or by durable,
high-value knowledge.

Production memory stacks answer this triage with one of two
single-factor policies. Retrieval-augmented systems rank stored items by
\emph{semantic similarity} to the current query
\citep{lewis2020rag,packer2023memgpt}; time-to-live and sliding-window
schemes rank by \emph{recency}. Both are intuitive and both are
mis-specified for the governing decision. Similarity-to-query is only
defined \emph{at retrieval time}, once the query exists --- but the
\emph{forgetting} decision is made earlier, at consolidation, when the
future query is unknown. Recency is query-agnostic but discards the
common case of an old-but-important fact (a user's allergy, a project
constraint) in favour of recent low-value chatter. Neither policy has a
notion of how \emph{useful} a memory is to future behaviour.

Human memory is not single-factor. Decades of work show that what
survives forgetting is governed by the \emph{value} of an item:
people preferentially retain information tagged as high-value even at
the expense of low-value detail \citep{castel2008value}; depth of
processing --- semantic, self-referential, and goal-relevant encoding
--- predicts durability over shallow perceptual encoding
\citep{craik1972lop,rogers1977sre}; emotional arousal modulates
consolidation \citep{mcgaugh2000consolidation}; and the very shape of
forgetting tracks the \emph{need-probability} of information in the
environment, an adaptive rather than passive decay
\citep{anderson1991adaptive,ebbinghaus1885forgetting}. The common thread
is that retention is driven by a multi-factor estimate of an item's
expected future usefulness, not by any single cue.

We bring this principle to agentic memory. We define a
\textbf{multi-factor memory value function}
\begin{equation}
V(m) \;=\; \sum_{i} w_i\, f_i(m),
\label{eq:value-intro}
\end{equation}
a weighted combination of seven interpretable factors --- emotional
intensity, goal relevance, value alignment, self/user relevance, task
utility, reliability, and usage history. A \emph{single} scalar $V(m)$
then uniformly controls all three triage operations: encoding depth,
forget risk, and retrieval rank. Crucially, the weights $w_i$ are not
hand-set: they are \textbf{learned} from downstream task return by a
gradient-free optimiser, because the
encode\,$\rightarrow$\,\allowbreak forget\,$\rightarrow$\,\allowbreak
retrieve\,$\rightarrow$\,\allowbreak answer
pipeline is non-differentiable. The value function is thus fit to
maximise task success, in the spirit of expected-utility credit
assignment \citep{sutton2018rl}.

Along the way we surface a methodological pitfall in how memory
retention is commonly measured, and turn it into a clean experimental
contrast. On the LongMemEval benchmark \citep{wu2025longmemeval}, the
``gold'' evidence for a question is, by construction, defined
\emph{relative to that question}. If a memory policy is allowed to score
goal relevance as the cosine similarity between a stored turn and the
held-out evaluation question --- an \emph{oracle} that peeks at the
future query --- then similarity alone retains essentially all gold
evidence ($\approx\!0.96$), and no multi-factor advantage is possible.
But this measures \emph{retrieval}, not forgetting: a real agent has not
seen the evaluation question at consolidation time. In the realistic
\emph{blind} regime, where goal relevance is computed only against
information available at consolidation (the topic of the ongoing
session), the picture inverts: similarity collapses to near-chance and a
learned multi-factor value retains gold far better than any single-factor
baseline. Separating these two regimes is, we argue, necessary for any
honest evaluation of a forgetting policy.

\paragraph{Contributions.}
\begin{enumerate}[leftmargin=1.4em,itemsep=2pt,topsep=2pt]
\item A \textbf{multi-factor memory value function}
(\cref{eq:value-intro}) for the expected usefulness of a memory to
future agent behaviour, with seven interpretable factors, and a single
scalar that uniformly drives encoding depth, forgetting, and retrieval
(\cref{sec:method}).
\item A \textbf{learned-weight} formulation (A2): a gradient-free
optimiser fits $\wvec$ to downstream task return, replacing hand-tuned
resistances with credit-assigned weights (\cref{sec:method:learn}).
\item A \textbf{methodological contrast} --- \emph{oracle} vs.\
\emph{blind} forgetting --- that separates retrieval from retention and
shows why query-defined retention benchmarks cannot, on their own,
demonstrate a forgetting advantage (\cref{sec:exp:blind}).
\item \textbf{Empirical results} on the real LongMemEval benchmark: in
the blind regime, learned multi-factor value retains
$0.811\pm0.047$ of gold evidence versus $0.633$ (uniform), $0.380$
(recency), and $\le0.39$ (any single factor), with interpretable learned
weights (\cref{sec:exp:blind}); a synthetic confound study validates the
learner (\cref{sec:exp:synth}).
\item An \textbf{open-source, CPU-only} substrate and evaluation harness;
every reported number reproduces with no API calls.
\end{enumerate}
